\begingroup
\footnotesize
\begin{thebibliography}{99}
\setlength{\itemsep}{1.5pt}

\bibitem{hart2000} S.~Hart, A.~Mas-Colell. A simple adaptive procedure leading to correlated equilibrium. \emph{Econometrica} 68(5), 2000.
\bibitem{zinkevich2007} M.~Zinkevich, M.~Johanson, M.~Bowling, C.~Piccione. Regret minimization in games with incomplete information. \emph{NIPS}, 2007.
\bibitem{neller2013} T.~Neller, M.~Lanctot. An introduction to counterfactual regret minimization. AAAI Model AI Assignments, 2013. \texttt{modelai.gettysburg.edu/2013/cfr}.
\bibitem{lanctot2009} M.~Lanctot, K.~Waugh, M.~Zinkevich, M.~Bowling. Monte Carlo sampling for regret minimization in extensive games. \emph{NIPS}, 2009.
\bibitem{tammelin2014} O.~Tammelin. Solving large imperfect information games using CFR+. arXiv:1407.5042, 2014.
\bibitem{tammelin2015} O.~Tammelin, N.~Burch, M.~Johanson, M.~Bowling. Solving heads-up limit Texas hold'em. \emph{IJCAI}, 2015.
\bibitem{bowling2015} M.~Bowling, N.~Burch, M.~Johanson, O.~Tammelin. Heads-up limit hold'em poker is solved. \emph{Science} 347(6218), 2015.
\bibitem{brown2019dcfr} N.~Brown, T.~Sandholm. Solving imperfect-information games via discounted regret minimization. \emph{AAAI}, 2019.
\bibitem{brown2015pruning} N.~Brown, T.~Sandholm. Regret-based pruning in extensive-form games. \emph{NIPS}, 2015.
\bibitem{johanson2012pcs} M.~Johanson, N.~Bard, M.~Lanctot, R.~Gibson, M.~Bowling. Efficient Nash equilibrium approximation through Monte Carlo counterfactual regret minimization / public chance sampling. \emph{AAMAS}, 2012.
\bibitem{koller1996} D.~Koller, N.~Megiddo, B.~von~Stengel. Efficient computation of equilibria for extensive two-person games. \emph{Games and Economic Behavior} 14, 1996.
\bibitem{schmid2019vr} M.~Schmid, N.~Burch, M.~Lanctot, M.~Moravcik, R.~Kadlec, M.~Bowling. Variance reduction in Monte Carlo CFR using baselines. \emph{AAAI}, 2019.
\bibitem{gilpin2007} A.~Gilpin, T.~Sandholm. Lossless abstraction of imperfect information games. \emph{Journal of the ACM} 54(5), 2007.
\bibitem{waugh2009} K.~Waugh, D.~Schnizlein, M.~Bowling, D.~Szafron. Abstraction pathologies in extensive games. \emph{AAMAS}, 2009.
\bibitem{waugh2013iso} K.~Waugh. A fast and optimal hand isomorphism algorithm. AAAI Workshop on Computer Poker, 2013.
\bibitem{johanson2013} M.~Johanson, N.~Burch, R.~Valenzano, M.~Bowling. Evaluating state-space abstractions in extensive-form games. \emph{AAMAS}, 2013.
\bibitem{ganzfried2013at} S.~Ganzfried, T.~Sandholm. Action translation in extensive-form games with large action spaces: axioms, paradoxes, and the pseudo-harmonic mapping. \emph{IJCAI}, 2013.
\bibitem{ganzfried2014} S.~Ganzfried, T.~Sandholm. Potential-aware imperfect-recall abstraction with earth mover's distance. \emph{AAAI}, 2014.
\bibitem{ganzfried2012purification} S.~Ganzfried, T.~Sandholm, K.~Waugh. Strategy purification and thresholding: effective non-equilibrium approaches for playing large games. \emph{AAMAS}, 2012.
\bibitem{brown2019deep} N.~Brown, A.~Lerer, S.~Gross, T.~Sandholm. Deep counterfactual regret minimization. \emph{ICML}, 2019.
\bibitem{steinberger2019} E.~Steinberger. Single deep counterfactual regret minimization. arXiv:1901.07621, 2019.
\bibitem{steinberger2020dream} E.~Steinberger, A.~Lerer, N.~Brown. DREAM: deep regret minimization with advantage baselines and model-free learning. arXiv:2006.10410, 2020.
\bibitem{mcaleer2022} S.~McAleer, G.~Farina, M.~Lanctot, T.~Sandholm. ESCHER: eschewing importance sampling in games by computing a history value function to estimate regret. \emph{ICLR}, 2023.
\bibitem{liu2022recfr} W.~Liu, B.~Li, J.~Togelius. Model-free neural counterfactual regret minimization with bootstrap learning. \emph{IEEE Transactions on Games}, 2022.
\bibitem{burch2014} N.~Burch, M.~Johanson, M.~Bowling. Solving imperfect-information games using decomposition. \emph{AAAI}, 2014.
\bibitem{deepstack} M.~Moravcik, M.~Schmid, N.~Burch, V.~Lisy, D.~Morrill, N.~Bard, T.~Davis, K.~Waugh, M.~Johanson, M.~Bowling. DeepStack: expert-level artificial intelligence in heads-up no-limit poker. \emph{Science} 356(6337), 2017.
\bibitem{libratus} N.~Brown, T.~Sandholm. Superhuman AI for heads-up no-limit poker: Libratus beats top professionals. \emph{Science} 359(6374), 2018. (Nested subgame solving: \emph{NIPS} 2017.)
\bibitem{brown2018dls} N.~Brown, T.~Sandholm, B.~Amos. Depth-limited solving for imperfect-information games. \emph{NeurIPS}, 2018.
\bibitem{pluribus} N.~Brown, T.~Sandholm. Superhuman AI for multiplayer poker. \emph{Science} 365(6456), 2019.
\bibitem{rebel} N.~Brown, A.~Bakhtin, A.~Lerer, Q.~Gong. Combining deep reinforcement learning and search for imperfect-information games. \emph{NeurIPS}, 2020.
\bibitem{sog} M.~Schmid et al. Student of Games: a unified learning algorithm for both perfect and imperfect information games. \emph{Science Advances} 9, 2023.
\bibitem{gtowizardai} GTO Wizard. GTO Wizard AI explained; How solvers work; How GTO Wizard solved PLO. Company engineering blog, 2023--2026 (vendor claims, no peer review). \texttt{blog.gtowizard.com}.
\bibitem{perolat2020} J.~Perolat, R.~Munos, J.-B.~Lespiau, et al. From Poincar\'e recurrence to convergence in imperfect information games: finding equilibrium via regularization. \emph{ICML}, 2021.
\bibitem{alphazero} D.~Silver et al. Mastering chess and shogi by self-play with a general reinforcement learning algorithm. arXiv:1712.01815, 2017; \emph{Science} 362, 2018.
\bibitem{nfsp} J.~Heinrich, D.~Silver. Deep reinforcement learning from self-play in imperfect-information games. arXiv:1603.01121, 2016.
\bibitem{psro} M.~Lanctot, V.~Zambaldi, A.~Gruslys, et al. A unified game-theoretic approach to multiagent reinforcement learning. \emph{NeurIPS}, 2017.
\bibitem{deepnash} J.~Perolat, B.~De~Vylder, D.~Hennes, et al. Mastering the game of Stratego with model-free multiagent reinforcement learning. \emph{Science} 378, 2022.
\bibitem{mmd} S.~Sokota, R.~D'Orazio, J.Z.~Kolter, N.~Loizou, M.~Lanctot, I.~Mitliagkas, N.~Brown, C.~Kroer. A unified approach to reinforcement learning, quantal response equilibria, and two-player zero-sum games. \emph{ICLR}, 2023.
\bibitem{alphaholdem} E.~Zhao, R.~Yan, J.~Li, K.~Li, J.~Xing. AlphaHoldem: high-performance artificial intelligence for heads-up no-limit poker via end-to-end reinforcement learning. \emph{AAAI}, 2022.
\bibitem{timbers2022} F.~Timbers, N.~Bard, E.~Lockhart, M.~Lanctot, M.~Schmid, N.~Burch, J.~Schrittwieser, T.~Hubert, M.~Bowling. Approximate exploitability: learning a best response in large games. \emph{IJCAI}, 2022.
\bibitem{vrpo} Z.~Fan, G.~Farina. GAE falls short in imperfect-information self-play reinforcement learning. arXiv preprint, 2026 (unreviewed).
\bibitem{johanson2011} M.~Johanson, K.~Waugh, M.~Bowling, M.~Zinkevich. Accelerating best response calculation in large extensive games. \emph{IJCAI}, 2011.
\bibitem{lbr} V.~Lisy, M.~Bowling. Equilibrium approximation quality of current no-limit poker bots (local best response). AAAI Workshop on Computer Poker, 2017.
\bibitem{aivat} N.~Burch, M.~Schmid, M.~Moravcik, M.~Bowling. AIVAT: a new variance reduction technique for agent evaluation in imperfect information games. \emph{AAAI}, 2018.
\bibitem{acpc} N.~Bard, J.~Hawkin, J.~Rubin, M.~Zinkevich. The annual computer poker competition. \emph{AI Magazine} 34(2), 2013.
\bibitem{ostroumov} O.~Ostroumov. How I made my first million: story of Draw Solver (2-7 triple draw / badugi). First-person engineering account, GipsyTeam/Medium, 2025 (practitioner source).
\bibitem{ho2015} K.~Ho. A no-limit Omaha hi-lo poker jam/fold endgame equilibrium. Master's thesis, Harvard Extension School, 2015.
\bibitem{pokercnn} N.~Yakovenko, L.~Cao, C.~Raffel, J.~Fan. Poker-CNN: a pattern learning strategy for making draws and bets in poker games. \emph{AAAI}, 2016.
\bibitem{zadeh} N.~Zadeh. Computation of optimal poker strategies. \emph{Operations Research} 25(4), 1977.
\bibitem{ginrummy} B.~Goldstein, J.-P.~Astudillo~Guerra, E.~Haigh, B.~Cruz~Ulloa, J.~Blum. Extracting learned discard and knocking strategies from a gin rummy bot. \emph{AAAI/EAAI}, 2021.
\bibitem{openspiel} M.~Lanctot, E.~Lockhart, et al. OpenSpiel: a framework for reinforcement learning in games. arXiv:1908.09453, 2019.
\bibitem{pokerrl} E.~Steinberger. PokerRL (with Deep-CFR/SD-CFR reference implementations). \texttt{github.com/EricSteinberger/PokerRL}, 2019.
\bibitem{rlcard} D.~Zha, K.-H.~Lai, Y.~Cao, S.~Huang, R.~Wei, J.~Guo, X.~Hu. RLCard: a toolkit for reinforcement learning in card games. AAAI-20 Workshop, 2019.
\bibitem{pokerkit} J.~Kim. PokerKit: a comprehensive Python library for fine-grained multi-variant poker game simulations. \emph{IEEE Transactions on Games}, 2023.
\bibitem{cactuskev} K.~Suffecool (Cactus Kev); P.~Senzee; and descendants (OMPEval, PokerHandEvaluator). Fast 5-card hand evaluation via 7{,}462 equivalence classes and perfect hashing. Community engineering lineage, 2006--. \texttt{suffe.cool/poker/evaluator.html}.

\end{thebibliography}
\endgroup
